\documentclass[10pt,a4paper]{amsart}
\usepackage[margin=19mm]{geometry}
\usepackage{amsmath,amssymb,mathtools}
\usepackage[scr=rsfs,cal=euler]{mathalfa}
\usepackage{xfrac}
\usepackage[unicode,hidelinks,bookmarks=false]{hyperref}
\newcommand{\ie}{i.e.\ }
\newcommand{\cf}{cf.\ }
\newcommand{\resp}{respectively\ }
\newcommand{\Subsection}{\subsection}
\providecommand{\nobreakdash}{\nobreak\hbox{-}}
\setlength{\emergencystretch}{2em}
\multlinegap0pt
\usepackage{amssymb}
\usepackage[shortlabels]{enumitem}
\usepackage{xcolor}
\usepackage{tikz}
\newtheorem{proposition}{Proposition}
\def\N{{\mathbb N}}
\def\R{{\mathbb R}}
\title{Pressure at infinity on countable Markov shifts}
\author{Anibal Velozo}
\date{Source: arXiv:2410.16552v1 (CC BY 4.0)}
\begin{document}
\maketitle

\noindent\emph{Source and licence.} Pressure at infinity on countable Markov shifts. Version arXiv:2410.16552v1 (CC BY 4.0), distributed by arXiv under the Creative Commons Attribution 4.0 International licence, http://creativecommons.org/licenses/by/4.0/, as recorded in the arXiv licence metadata for this exact version and shown on the abstract page https://arxiv.org/abs/2410.16552v1. Full licence notice: \texttt{licenses/arxiv-2410.16552-CC-BY-4.0.txt}.\par\smallskip

\noindent\textit{Editorial scope.} Counted unit: the article's proposition inf (source0.tex lines 369--370). The article's complete original proof of it is delivered with it (source0.tex lines 371--391, one \texttt{proof} environment of 21 lines). No mathematics is written, repaired or re-derived: the statement and the proof are the article's own lines, in the article's own notation, and no retained line is substituted or re-flowed.  No further source-local context is needed: every cross-reference of every retained line resolves inside the two delivered spans.  Nothing was omitted from any retained span, so the package carries no omission record.  No retained line contains a citation, so the wrapper carries no bibliography and no bibliography entry.  No retained line contains a figure environment, an image-inclusion command or drawing code, so no image is delivered and the package declares no asset dependency.  Editorial formatting only, in addition: an AMS \texttt{amsart} preamble that restates the article's own declarations and macros used by the retained lines, namely the environments \texttt{proposition} on separate counters, together with the standard packages those lines need.  The retained environments print in this document as Proposition 1 in order of appearance, whereas the article prints the same items with the numbering of its own full document.  The complete native source of the article as distributed in the arXiv e-print for this exact version is bundled unchanged as \texttt{sources/167206/source0.tex}.
\begin{proposition}\label{inf} Let $\phi:\Sigma\to\R$ be a potential with summable variations. Let $V=\sum_{k=1}^\infty \frac{1}{k} 1_{[k]}\in C_0(\Sigma)$. Then, $$P^{top}_\infty(\phi,a)=\lim_{t\to\infty} P(\phi-tV).$$
\end{proposition}

\begin{proof} We will first prove that $P^{top}_\infty(\phi,a)=P^{top}_\infty(\phi-tV,a)$, for every $t\ge 0$. Observe that if $x\in \text{Per}_a(q,M,n)$, then $S_nV(x)\le \frac{n}{M}+\frac{n}{q}$. It follows that,
$$Z_n(\phi-tV,a;q,M)\ge Z_n(\phi,a;q,M)\exp\bigg(-nt\bigg(\frac{1}{M}+\frac{1}{q}\bigg)\bigg),$$
and therefore,
$$P^{top}_\infty(\phi-tV,a;q,M)\ge P^{top}_\infty(\phi,a;q,M)-t\bigg(\frac{1}{M}+\frac{1}{q}\bigg).$$
Sending $q$ and $M$ to infinity we obtain $P^{top}_\infty(\phi-tV,a)\ge P^{top}_\infty(\phi,a)$. The other inequality follows from $\phi\ge \phi-tV$. We conclude that $P_\infty^{top}(\phi,a)=P_\infty^{top}(\phi-tV,a)$. 

Since $P^{top}_\infty(\phi,a)=P^{top}_\infty(\phi-tV,a)\le P(\phi-tV)$, we have $P^{top}_\infty(\phi,a)\le \lim_{t\to\infty}P(\phi-tV)$. In order to complete the proof, we need to show that  $\lim_{t\to\infty}P(\phi-tV)\le P^{top}_\infty(\phi,a)$. 

Fix $q,M\in\N$. If $x\in \text{Per}_a^c(q,M,n)$, then $\#\{1\le k\le n:x_k\le q\}> \frac{n}{M}$. In particular, $S_nV(x)> \frac{n}{Mq}.$ Then,
$$\sum_{x\in \text{Per}_a^c(q,M,n)}\exp(S_n(\phi-tV)(x))\le \exp\bigg(-\frac{tn}{Mq}\bigg)\sum_{x\in \text{Per}_a^c(q,M,n)}\exp(S_n\phi(x)).$$
%$$\sum_{x\in Per^c(q,M,n)}\exp(F_n(x)-tV_n(x))\le \#Per^c(q,M,n)\exp(n(||F||_0-\frac{t}{Mq})).$$
By definition of the pressure, there exists $n_0\in\N$ such that if $n\ge n_0,$ then \\$\sum_{x\in \text{Per}_a^c(q,M,n)}\exp(S_n\phi(x))\le \exp(n(P(\phi)+1)).$ Therefore,
%$$\#Per^c(q,M,n)\le \exp(n(h_{top}(\Sigma,\sigma)+1)).$$ In particular 
\begin{align*}\sum_{x\in \text{Per}_a^c(q,M,n)}\exp(S_n(\phi-tV)(x))\le \exp\bigg(n\bigg(P(\phi)+1-\frac{t}{Mq}\bigg)\bigg),\end{align*}
%for every $n\ge n_0$. Choose $t=t(q,M)>(Mq)(h_{top}(\Sigma,\sigma)+2+||F||_0)$. In this situation we have that
for every $n\ge n_0$. Finally, for $n\ge n_0$ and $t>qM(P(\phi)+1+\kappa)$, we have
\begin{align*} Z_n(\phi-tV,a)&=\sum_{x\in \text{Per}_a(q,M,n)}\exp(S_n(\phi-tV)(x))+\sum_{x\in \text{Per}_a^c(q,M,n)}\exp(S_n(\phi-tV)(x))\\
&\le \sum_{x\in \text{Per}_a(q,M,n)}\exp(S_n\phi(x))+e^{-\kappa n},
 \end{align*}
for every $\kappa>0$. We conclude that $P(\phi-tV)\le \max\{P^{top}_\infty(\phi,a;q,M),-\kappa\}$, and therefore  $\lim_{t\to\infty}P(\phi-tV)\le P^{top}_\infty(\phi,a;q,M)$. 
\end{proof}

\end{document}
